%%
%% Design Considerations for a Full Numeric Tower in MScheme
%%
\documentclass[12pt]{report}
\usepackage[margin=1in]{geometry}
\usepackage{longtable}
\usepackage{booktabs}
\usepackage{listings}
\usepackage{xcolor}
\usepackage{hyperref}
\usepackage{textcomp}
\usepackage{amssymb}
\usepackage{amsmath}
\usepackage{enumitem}

\lstdefinelanguage{Scheme}{
  morekeywords={define,lambda,if,cond,else,let,let*,letrec,begin,
    set!,quote,quasiquote,unquote,unquote-splicing,and,or,not,
    case,do,delay,force,map,for-each,apply,call-with-current-continuation,
    call/cc,require-modules,load,macro,define-global-symbol,
    exact->inexact,inexact->exact,make-rectangular,make-polar,
    real-part,imag-part,magnitude,angle,numerator,denominator,
    rationalize,number->string,string->number},
  sensitive=true,
  morecomment=[l]{;},
  morestring=[b]",
  literate={lambda}{$\lambda$}1,
}

\lstdefinelanguage{Modula3}{
  morekeywords={MODULE,INTERFACE,IMPORT,FROM,PROCEDURE,BEGIN,END,
    VAR,CONST,TYPE,RECORD,OBJECT,METHODS,OVERRIDES,RETURN,
    IF,THEN,ELSE,ELSIF,WHILE,DO,FOR,TO,BY,LOOP,EXIT,REPEAT,UNTIL,
    CASE,OF,WITH,TRY,EXCEPT,FINALLY,RAISE,RAISES,
    NEW,ARRAY,REF,REFANY,SET,BOOLEAN,INTEGER,REAL,LONGREAL,EXTENDED,
    CHAR,TEXT,LONGINT,
    TRUE,FALSE,NIL,BRANDED,REVEAL,EVAL,NOT,AND,OR,IN,DIV,MOD,
    ISTYPE,NARROW,TYPECODE,GENERIC,ADDRESS,BITS,BITSIZE},
  sensitive=true,
  morecomment=[n]{(*}{*)},
  morestring=[b]",
}

\lstset{
  basicstyle=\ttfamily\small,
  breaklines=true,
  frame=single,
  backgroundcolor=\color{gray!5},
  keywordstyle=\bfseries\color{blue!70!black},
  commentstyle=\itshape\color{green!50!black},
  stringstyle=\color{red!60!black},
  showstringspaces=false,
  columns=flexible,
  tabsize=2,
  xleftmargin=2em,
  framexleftmargin=1.5em,
}

\title{\bfseries\sffamily\fontsize{24.88}{30}\selectfont
  Design Considerations for a\\Full Numeric Tower in MScheme}
\author{Mika Nystr\"om \\ {\tt mika@alum.mit.edu}}
\date{March 2026}

\begin{document}
\maketitle
\thispagestyle{empty}

\newpage
\tableofcontents
\newpage

%======================================================================
\chapter{Introduction}
%======================================================================

MScheme is a Scheme interpreter and compiler written in Modula-3,
descended from Peter Norvig's JScheme~\cite{jscheme}.  Its current
numeric representation is a single type: IEEE~754 double precision
(Modula-3 {\tt LONGREAL}), boxed as {\tt REF LONGREAL} and stored
in the universal {\tt SchemeObject.T = REFANY} slot.  This
representation is simple and fast but falls short of
the numeric tower specified by R5RS~\cite{r5rs} and later standards.

This report examines the design space for extending MScheme with a
full numeric tower---exact integers of arbitrary precision, exact
rationals, inexact reals, and complex numbers---while preserving
the system's defining characteristics:

\begin{enumerate}
\item \emph{Modula-3 interoperability.}  Every Scheme value is a
  {\tt REFANY}, passable to and from Modula-3 code via the sstubgen
  stub generator without marshalling.

\item \emph{Compiled-code performance.}  The MScheme compiler
  performs type inference and LONGREAL unboxing in self-tail-call loops,
  yielding arithmetic throughput competitive with native Modula-3.
  A numeric tower must not destroy this property.

\item \emph{Simplicity.}  MScheme's implementation is compact
  (approximately 5\,000 lines of Modula-3 for the interpreter,
  3\,000 lines of Scheme for the compiler).  Complexity must be
  proportional to value delivered.
\end{enumerate}

The report surveys the relevant standards (Chapter~\ref{ch:standards}),
analyzes representation choices (Chapter~\ref{ch:representations}),
examines Modula-3 interoperability constraints
(Chapter~\ref{ch:m3interop}), considers compiler implications
(Chapter~\ref{ch:compiler}), develops a concrete integer
representation design (Chapter~\ref{ch:integer-design}), briefly
reviews exotic representations (Chapter~\ref{ch:exotic}), discusses
verification and testing (Chapter~\ref{ch:testing}), and concludes
with recommendations (Chapter~\ref{ch:conclusions}).
A companion document, \emph{Numeric Tower Phase~1: Implementation
Proposal}, provides the concrete code-level plan for implementing
exact integers.

%======================================================================
\chapter{Standards Compliance}
\label{ch:standards}
%======================================================================

%----------------------------------------------------------------------
\section{R4RS}
%----------------------------------------------------------------------

The Revised$^4$ Report~\cite{r4rs} permits considerable latitude in
the numeric tower.  Section~6.5 states:

\begin{quote}
  It is not necessary for [an implementation] to implement the
  whole tower of subtypes \ldots\ These types form a tower of
  subtypes in which each level is a subset of the level above it:
  number $\supset$ complex $\supset$ real $\supset$ rational $\supset$
  integer.
\end{quote}

R4RS requires that {\tt number?}, {\tt complex?}, {\tt real?},
{\tt rational?}, and {\tt integer?} be provided, and that {\tt exact?}
and {\tt inexact?} classify numbers.  Procedures that operate on
missing types (e.g., {\tt numerator}/{\tt denominator} without
rationals) are explicitly non-essential.

MScheme's current single-{\tt LONGREAL} representation satisfies R4RS
at a minimal level: all required predicates exist, arithmetic is
complete, but the exact/inexact distinction is collapsed.  A companion
report~\cite{r4rsconf} documents the full audit.

%----------------------------------------------------------------------
\section{R5RS}
%----------------------------------------------------------------------

The Revised$^5$ Report~\cite{r5rs} tightens the requirements.  Key
changes relevant to the numeric tower:

\begin{itemize}
\item {\tt eqv?} on numbers must distinguish exact from inexact:
  {\tt (eqv?\ 1\ 1.0)} must return {\tt \#f} (R5RS~\S6.1).
  MScheme currently returns {\tt \#t}.

\item {\tt exact->inexact} and {\tt inexact->exact} are required
  (not merely permitted).  MScheme now provides them as identity
  functions, which is incorrect once distinct exact types exist.

\item The specification states that ``an implementation \emph{may}
  use fixnums, flonums, or perhaps some other representation for
  numbers, as long as the semantic requirements are met.''  Crucially,
  exact arithmetic on integers must not lose precision:
  {\tt (+\ \mbox{MOST-POSITIVE-FIXNUM}\ 1)} must produce a correct
  result, not wrap around or signal an error.

\item {\tt number->string} and {\tt string->number} must respect
  exactness.  A string like {\tt "1/3"} should parse to an exact
  rational if the implementation supports them.
\end{itemize}

%----------------------------------------------------------------------
\section{R7RS}
%----------------------------------------------------------------------

R7RS-small~\cite{r7rs} largely inherits the R5RS numeric tower but
adds explicit library structure.  The base library {\tt (scheme base)}
requires exact integers; {\tt (scheme complex)} provides
{\tt make-rectangular}, {\tt make-polar}, etc.; {\tt (scheme inexact)}
provides transcendentals.  This modular structure maps naturally to a
phased implementation: exact integers first, then optionally rationals
and complex numbers.

%----------------------------------------------------------------------
\section{Relevant SRFIs}
%----------------------------------------------------------------------

Several SRFIs bear on numeric tower design:

\begin{description}[style=nextline]
\item[SRFI-60: Integers as Bits]
  Bitwise operations on exact integers.
  Requires an exact integer type.

\item[SRFI-151: Bitwise Operations]
  Supersedes SRFI-60 with a cleaner API.  Same requirement.

\item[SRFI-70: Numbers]
  Proposes that inexact numbers include $+\infty$, $-\infty$, and NaN
  as first-class values.  Already satisfied by IEEE~754 doubles.

\item[SRFI-77/SRFI-143: Fixnums]
  Defines a fixnum-specific API.  Useful once the implementation
  distinguishes fixnums from bignums.

\item[SRFI-144: Flonums]
  Defines a flonum-specific API over IEEE~754 doubles.
  Compatible with the current representation.
\end{description}

%======================================================================
\chapter{Representation Choices}
\label{ch:representations}
%======================================================================

%----------------------------------------------------------------------
\section{Current State: {\tt REF LONGREAL} Only}
%----------------------------------------------------------------------

Every Scheme number is currently a {\tt REF LONGREAL}, allocated on
the Modula-3 traced heap.  The interface {\tt SchemeLongReal}
provides:

\begin{lstlisting}[language=Modula3]
TYPE T = REF LONGREAL;

PROCEDURE FromLR(x : LONGREAL) : T;
PROCEDURE FromI(x : INTEGER) : T;
PROCEDURE FromO(x : SchemeObject.T) : LONGREAL RAISES { E };
\end{lstlisting}

\noindent {\tt FromLR} caches singletons for $-1$, $0$, $1$, $2$ to
reduce allocation pressure.  {\tt FromO} performs
{\tt NARROW(x, T)\^{}} to extract the {\tt LONGREAL} value, serving
as the central bridge from the polymorphic Scheme world to typed
Modula-3 arithmetic.

The {\tt IsExact} predicate in {\tt SchemePrimitive.m3} tests whether a
double is integer-valued and within the IEEE~754 exact-integer range
($|x| < 2^{53}$):

\begin{lstlisting}[language=Modula3]
PROCEDURE IsExact(x : Object) : BOOLEAN RAISES { E } =
  BEGIN
    ...
    WITH d = FromO(x) DO
      RETURN d = FLOAT(ROUND(d),LONGREAL) AND
             ABS(d) < 102962884861573423.0d0
    END
  END IsExact;
\end{lstlisting}

\noindent This conflates ``integer-valued'' with ``exact,'' which is
correct only in the absence of distinct exact types.

%----------------------------------------------------------------------
\section{Exact Integers}
\label{sec:exact-integers}
%----------------------------------------------------------------------

\subsection{Fixnums: {\tt REF INTEGER}}

The simplest exact integer representation is {\tt REF INTEGER}.
On 64-bit platforms (all current CM3 targets), {\tt INTEGER} is 64
bits, giving an exact range of $[-2^{63}, 2^{63}-1]$.  This
representation:

\begin{itemize}
\item Is a {\tt REFANY}---satisfies the interoperability invariant.
\item Is trivially distinguishable from {\tt REF LONGREAL}
  via {\tt ISTYPE}.
\item Supports native arithmetic ({\tt +}, {\tt -}, {\tt *},
  {\tt DIV}, {\tt MOD}) at full hardware speed.
\item Requires overflow detection to promote to bignums.
\end{itemize}

The overflow detection problem is non-trivial in Modula-3.  The
language does not expose carry or overflow flags; the only portable
approach is pre-check or post-check using {\tt Word.T} arithmetic.
Section~\ref{sec:overflow} discusses this in detail.

\subsection{Bignums: GMP ({\tt Mpz.T}) vs.\ Pure Modula-3 ({\tt BigInt.T})}

Two arbitrary-precision integer libraries exist in the CM3 ecosystem.

\paragraph{Mpz.T (GMP bindings).}
Located at {\tt intel-async/\allowbreak async-toolkit/\allowbreak
m3utils/\allowbreak m3utils/\allowbreak gmp/\allowbreak mpz/}.  The type
{\tt Mpz.T <: REFANY} wraps GMP's {\tt mpz\_t} with Modula-3 traced
references and weak-reference-based deallocation.  The bindings are
\emph{auto-generated} from a Scheme-based description of the GMP API
({\tt mpz.scm} $\to$ {\tt make-mpz.scm}), which means they track GMP
updates by regeneration rather than manual editing.

The interface follows GMP's output-first convention:

\begin{lstlisting}[language=Modula3]
TYPE T <: REFANY;

PROCEDURE NewInt(int : INTEGER) : T;
PROCEDURE add(f0, f1, f2 : T);   (* f0 := f1 + f2 *)
PROCEDURE mul(f0, f1, f2 : T);   (* f0 := f1 * f2 *)
PROCEDURE sub(f0, f1, f2 : T);   (* f0 := f1 - f2 *)
PROCEDURE tdiv_q(f0, f1, f2 : T); (* truncated quotient *)
PROCEDURE ToInteger(t : T) : INTEGER;
PROCEDURE get_d(t : T) : LONGREAL;  (* float approx *)
PROCEDURE FormatDecimal(t : T) : TEXT;
\end{lstlisting}

\noindent The full API exposes over 200 GMP operations, organized
into categories:

\begin{description}[style=nextline,leftmargin=1.5em]
\item[Arithmetic.]
  {\tt add}, {\tt sub}, {\tt mul}, {\tt neg}, {\tt abs},
  plus {\tt add\_ui}/{\tt mul\_ui} variants for mixed-precision
  operations, and {\tt addmul}/{\tt submul} fused operations.

\item[Division.]
  Three rounding modes: ceiling ({\tt cdiv\_*}), floor ({\tt fdiv\_*}),
  and truncation ({\tt tdiv\_*}), each with quotient, remainder, and
  combined forms.  Plus {\tt divexact} for known-exact divisions
  and power-of-two variants ({\tt *\_2exp}).

\item[Bitwise.]
  {\tt and}, {\tt ior}, {\tt xor}, {\tt com} (complement),
  {\tt setbit}, {\tt clrbit}, {\tt combit}, {\tt tstbit},
  {\tt scan0}, {\tt scan1}, {\tt popcount}, {\tt hamdist}.
  Left/right shift via {\tt mul\_2exp} and {\tt fdiv\_q\_2exp}.
  These directly support SRFI-151 bitwise operations on exact
  integers.

\item[Number theory.]
  {\tt gcd}, {\tt gcdext} (extended Euclidean), {\tt lcm},
  {\tt invert} (modular inverse), {\tt powm} (modular
  exponentiation), {\tt jacobi}/{\tt kronecker},
  {\tt probab\_prime\_p}/{\tt millerrabin} (primality),
  {\tt nextprime}, {\tt remove} (factor extraction).

\item[Combinatorial.]
  {\tt fac\_ui} (factorial), {\tt bin\_ui}/{\tt bin\_uiui}
  (binomial coefficients), {\tt fib\_ui}/{\tt fib2\_ui}
  (Fibonacci), {\tt lucnum\_ui} (Lucas numbers).

\item[Conversion.]
  {\tt get\_si}/{\tt get\_ui} (to machine integer),
  {\tt get\_d} (to {\tt LONGREAL}), {\tt init\_set\_d}
  (from {\tt LONGREAL}), {\tt FormatDecimal}/{\tt
  FormatHexadecimal}/{\tt FormatBased} (to text in any base
  2--62), {\tt InitScan} (from text).
  {\tt ToInteger} and {\tt ToWord} are checked wrappers that
  raise a runtime error if the value is out of range.

\item[Size tests.]
  {\tt fits\_slong\_p}, {\tt fits\_uint\_p}, etc.---useful for
  the fixnum demotion check in the three-level design.
\end{description}

\paragraph{Memory layout.}
The internal representation ({\tt MpzRep.i3}) embeds GMP's
{\tt mpz\_t} structure directly in the Modula-3 object:

\begin{lstlisting}[language=Modula3]
REVEAL Mpz.T = BRANDED "Mpz" OBJECT
    val : Val;
  END;
TYPE Val = RECORD
    w : ARRAY [0..2] OF Word.T;  (* >= sizeof(mpz_t) *)
  END;
\end{lstlisting}

\noindent On 64-bit platforms, each {\tt Mpz.T} occupies approximately
40--48~bytes (Modula-3 object header plus three machine words for the
{\tt mpz\_t}), plus whatever GMP allocates internally for limbs beyond
the inline storage.  Cleanup is automatic: every {\tt Mpz.T} is
registered with a {\tt WeakRef} callback that calls {\tt
\_\_gmpz\_clear} when the object is collected.

GMP is among the fastest bignum libraries available.
The external dependency is the main drawback.

\paragraph{Mpq: rational bindings.}
The auto-generator's database ({\tt mpz.scm}) contains entries for
GMP's {\tt mpq} rational functions ({\tt mpq\_add}, {\tt mpq\_mul},
{\tt mpq\_canonicalize}, {\tt mpq\_get\_num}/{\tt mpq\_get\_den},
etc.), but these are \textbf{not currently exposed} in the Modula-3
interfaces.  Creating {\tt Mpq.T} bindings would be straightforward:
regenerate from the existing data, define {\tt Mpq.T <: REFANY}
analogously to {\tt Mpz.T}, and add {\tt WeakRef} cleanup.  This
would provide an alternative to {\tt BigRational} with GMP's
optimized rational arithmetic (automatic GCD reduction, efficient
numerator/denominator extraction).

\paragraph{MPFR: arbitrary-precision floats.}
MPFR bindings also exist ({\tt m3utils/mpfr/} and {\tt
caltech-other/mpfr/}), providing {\tt Mpfr.T <: REFANY} with
configurable precision, seven rounding modes, and the standard
transcendental functions.  While not needed for the core numeric
tower, MPFR could serve as the foundation for an optional
``exact real'' or ``arbitrary-precision float'' extension.

\paragraph{BigInt.T (pure Modula-3).}
Located at {\tt caltech-parser/\allowbreak cit\_util/\allowbreak src/}.  The
type {\tt BigInt.T <: ROOT} uses a sign-magnitude representation with
a dynamically-sized array of machine words:

\begin{lstlisting}[language=Modula3]
REVEAL BigInt.T = BRANDED "BigInt" OBJECT
    sign : [-1 .. 1];
    rep  : NSeq;
  END;
\end{lstlisting}

\noindent BigInt provides the basic operations ({\tt Add}, {\tt Sub},
{\tt Mul}, {\tt Div}, {\tt Mod}) and bitwise operations.  It is pure
Modula-3 with no external dependencies, making it highly portable.
However, it lacks the algorithmic optimizations of GMP (sub-quadratic
multiplication, etc.) and has seen less testing.

\paragraph{Recommendation.}
Use {\tt Mpz.T} as the primary bignum representation, with {\tt BigInt.T}
available as a fallback for platforms without GMP.  The {\tt Mpz.T
<: REFANY} typing is directly compatible with {\tt SchemeObject.T};
{\tt BigInt.T <: ROOT <: REFANY} is equally compatible.  Both can be
tested with {\tt ISTYPE} at dispatch points.

%----------------------------------------------------------------------
\section{Inexact Reals: {\tt LONGREAL}'s New Role}
%----------------------------------------------------------------------

With exact integers handled by {\tt REF INTEGER} and bignums,
{\tt REF LONGREAL} becomes the \emph{inexact} real type only.
A {\tt LONGREAL} value is never ``exact,'' even when integer-valued.
This is a semantic shift: currently {\tt (exact?\ 1.0)} returns
{\tt \#t}, but after the tower is implemented it must return {\tt \#f}
because the value {\tt 1.0} entered the system as inexact.

The {\tt SchemeLongReal.FromO} function remains the critical bridge
for Modula-3 interoperability.  It must be extended to handle exact
integer types:

\begin{lstlisting}[language=Modula3]
PROCEDURE FromO(x : SchemeObject.T) : LONGREAL RAISES { E } =
  BEGIN
    IF ISTYPE(x, SchemeLongReal.T) THEN
      RETURN NARROW(x, SchemeLongReal.T)^
    ELSIF ISTYPE(x, REF INTEGER) THEN
      RETURN FLOAT(NARROW(x, REF INTEGER)^, LONGREAL)
    ELSIF ISTYPE(x, Mpz.T) THEN
      RETURN Mpz.GetD(x)  (* GMP double extraction *)
    ELSE
      RAISE E(...)
    END
  END FromO;
\end{lstlisting}

\noindent This is the backward-compatibility linchpin: all existing
Modula-3 code that calls {\tt FromO} continues to work without
changes, receiving the {\tt LONGREAL} value regardless of the
underlying exact type.

%----------------------------------------------------------------------
\section{Exact Rationals}
%----------------------------------------------------------------------

\subsection{The Rational Generic}

The {\tt cit\_util} library provides {\tt Rational.ig}, a generic
rational-number interface parameterized over a base integer type:

\begin{lstlisting}[language=Modula3]
GENERIC INTERFACE Rational(BaseInt);
TYPE T = RECORD
    n : BaseInt.T;        (* numerator *)
    d : BaseInt.Natural;  (* denominator, always > 0 *)
  END;
\end{lstlisting}

\noindent Two instantiations exist:
\begin{itemize}
\item {\tt IntRational = Rational(IntForRat)} --- uses machine
  {\tt INTEGER}.
\item {\tt BigRational = Rational(BigInt)} --- uses arbitrary-precision
  {\tt BigInt.T}.
\end{itemize}

\subsection{The IntForRat Overflow Bug}
\label{sec:intforrat-bug}

{\tt IntForRat} adapts {\tt INTEGER} for use with the {\tt Rational}
generic.  Its interface contains an explicit warning:

\begin{quote}
  ``This implementation truncates its results MOD NUMBER(INTEGER) so
  that they are between FIRST(INTEGER) and LAST(INTEGER) inclusive.''
\end{quote}

\noindent The arithmetic operations are direct wrappers:

\begin{lstlisting}[language=Modula3]
PROCEDURE Mul(a, b : T) : T = BEGIN RETURN a * b END Mul;
PROCEDURE Add(a, b : T) : T = BEGIN RETURN a + b END Add;
\end{lstlisting}

\noindent These \emph{silently overflow}.  When {\tt Rational.Add}
calls {\tt IntForRat.Mul} to compute cross-products of numerators
and denominators, intermediate values can wrap around $2^{64}$,
producing silently wrong results.  {\tt IntRational} is therefore
\textbf{unusable} for a correct numeric tower.

\subsection{Recommendation}

Use {\tt BigRational} (instantiated with {\tt BigInt}) or create a
new instantiation {\tt MpzRational = Rational(MpzForRat)} where
{\tt MpzForRat} adapts {\tt Mpz.T} to the {\tt Rational} generic's
interface.  Since rational denominators are always positive, the
GCD-based reduction in {\tt Rational.mg} automatically maintains
lowest terms.

For Scheme integration, exact rationals should be boxed as {\tt REF
BigRational.T} (or a branded object type), distinguishable by
{\tt ISTYPE} at dispatch points.

%----------------------------------------------------------------------
\section{Complex Numbers}
%----------------------------------------------------------------------

Complex numbers are the least urgent component of the tower.
R5RS makes {\tt make-rectangular}, {\tt make-polar},
{\tt real-part}, {\tt imag-part}, {\tt magnitude}, and {\tt angle}
non-essential when an implementation does not support complex numbers.

When the time comes, a natural representation is:

\begin{lstlisting}[language=Modula3]
TYPE SchemeComplex = BRANDED "SchemeComplex" REF RECORD
    re, im : SchemeObject.T;
  END;
\end{lstlisting}

\noindent where {\tt re} and {\tt im} are themselves Scheme numbers
(exact or inexact).  This allows exact Gaussian integers
(e.g., $3+4i$ with exact components) alongside inexact complex
values.  The exact/inexact classification of a complex number is
determined by its components: exact iff both parts are exact.

%----------------------------------------------------------------------
\section{{\tt FLOAT} and {\tt EXTENDED}}
\label{sec:float-extended}
%----------------------------------------------------------------------

Modula-3 defines three floating-point types:
\begin{itemize}
\item {\tt REAL} --- IEEE~754 single precision (32-bit).
\item {\tt LONGREAL} --- IEEE~754 double precision (64-bit).
\item {\tt EXTENDED} --- ``extended precision,'' platform-dependent.
\end{itemize}

On most modern platforms (AMD64, ARM64), {\tt EXTENDED} is mapped
to {\tt LONGREAL} (both 64-bit).  Only on x86 (I386) targets does
{\tt EXTENDED} correspond to the 80-bit x87 extended format.

For the numeric tower, {\tt LONGREAL} is the only relevant inexact
type.  {\tt REAL} offers no advantage (less precision, same allocation
overhead), and {\tt EXTENDED}'s platform-dependent behavior makes it
unsuitable for portable Scheme semantics.  However, {\tt FromO} could
be extended to accept {\tt REF REAL} and {\tt REF EXTENDED} values
passed in from Modula-3 code, converting them to {\tt LONGREAL} for
Scheme consumption.

%======================================================================
\chapter{Modula-3 Interoperability}
\label{ch:m3interop}
%======================================================================

%----------------------------------------------------------------------
\section{The {\tt REFANY} Principle}
%----------------------------------------------------------------------

MScheme's defining interoperability feature is that every Scheme value
is a Modula-3 {\tt REFANY}.  This means:

\begin{itemize}
\item Scheme values can be passed to Modula-3 procedures expecting
  {\tt REFANY} (or specific reference types) without marshalling.
\item Modula-3 values that are reference types can be passed into
  Scheme without wrapping.
\item The sstubgen stub generator can produce Scheme bindings for
  Modula-3 interfaces automatically, handling numeric conversions
  via {\tt SchemeLongReal.FromO}, {\tt SchemeLongReal.Int}, etc.
\end{itemize}

Any numeric tower extension must preserve this principle.  Since
{\tt REF INTEGER <: REFANY}, {\tt Mpz.T <: REFANY}, and
{\tt BigInt.T <: ROOT <: REFANY}, all proposed exact types satisfy
the constraint.

%----------------------------------------------------------------------
\section{The sstubgen Bridge}
%----------------------------------------------------------------------

The sstubgen stub generator produces Modula-3 procedures that convert
between Scheme objects and typed Modula-3 values.  For numeric parameters,
it calls:

\begin{itemize}
\item {\tt SchemeLongReal.FromO(x)} for {\tt LONGREAL} parameters.
\item {\tt SchemeLongReal.Int(x)} for {\tt INTEGER} parameters.
\item {\tt SchemeLongReal.Card(x)} for {\tt CARDINAL} parameters.
\end{itemize}

With a numeric tower, these functions must accept the new types.
{\tt FromO} becomes the universal numeric extraction point:
regardless of whether a Scheme number is a fixnum, a bignum, or a
float, {\tt FromO} returns a {\tt LONGREAL} approximation.
{\tt Int} and {\tt Card} should preferentially extract exact integers
without the lossy round-trip through {\tt LONGREAL}:

\begin{lstlisting}[language=Modula3]
PROCEDURE Int(x : SchemeObject.T) : INTEGER RAISES { E } =
  BEGIN
    IF ISTYPE(x, REF INTEGER) THEN
      RETURN NARROW(x, REF INTEGER)^
    ELSIF ISTYPE(x, Mpz.T) THEN
      RETURN Mpz.ToInteger(x)  (* raises if out of range *)
    ELSIF ISTYPE(x, SchemeLongReal.T) THEN
      (* existing LONGREAL path *)
      ...
    END
  END Int;
\end{lstlisting}

%----------------------------------------------------------------------
\section{The {\tt LONGINT} Problem}
%----------------------------------------------------------------------

Modula-3 defines {\tt LONGINT} as a 64-bit integer type, distinct
from {\tt INTEGER} (which is also 64-bit on 64-bit targets).  On
32-bit targets, {\tt INTEGER} is 32 bits while {\tt LONGINT} remains
64 bits.

Currently, sstubgen does not handle {\tt LONGINT} parameters.  With
a numeric tower, a {\tt LONGINT}-accepting stub could extract a 64-bit
exact integer directly from a fixnum without going through {\tt LONGREAL}
(which loses precision beyond $2^{53}$).  This is a secondary concern
since 32-bit targets are increasingly rare, but the design should
not preclude it.

%----------------------------------------------------------------------
\section{{\tt Word.T} and Unsigned Arithmetic}
%----------------------------------------------------------------------

Modula-3's {\tt Word.T} is the unsigned machine word type, equal in
size to {\tt INTEGER} but with unsigned semantics.  The {\tt Word}
interface provides modular arithmetic ({\tt Word.Plus}, {\tt Word.Times},
{\tt Word.And}, etc.) that wraps on overflow.

Scheme has no native unsigned integer type, but SRFI-60 and SRFI-151
define bitwise operations on exact integers that assume a two's
complement representation extended to arbitrary precision.  The
implementation can use {\tt Word.T} arithmetic internally for
performance-critical bitwise operations while maintaining the Scheme
semantics of infinite-precision two's complement.

%======================================================================
\chapter{Compiler Implications}
\label{ch:compiler}
%======================================================================

%----------------------------------------------------------------------
\section{Current Architecture}
%----------------------------------------------------------------------

The MScheme compiler ({\tt schemecompiler/src/compiler.scm}, $\approx
3{,}000$ lines) translates Scheme source to Modula-3 source, which is
then compiled by the CM3 compiler.  All Scheme values are {\tt
SchemeObject.T = REFANY} at the Modula-3 level.

The compiler performs two key optimizations for numeric code:

\paragraph{Inline arithmetic.}
Binary operations {\tt +}, {\tt -}, {\tt *}, {\tt quotient},
{\tt remainder} and comparisons {\tt =}, {\tt <}, {\tt >}, {\tt <=},
{\tt >=} are compiled inline when applied to exactly two arguments.
The generated code extracts {\tt LONGREAL} values via {\tt NARROW},
performs native arithmetic, and re-boxes the result:

\begin{lstlisting}[language=Modula3]
(* compiled from (+ a b) *)
SchemeLongReal.FromLR(
  NARROW(a_0, SchemeLongReal.T)^ +
  NARROW(a_1, SchemeLongReal.T)^)
\end{lstlisting}

\paragraph{LONGREAL unboxing.}
For self-tail-call procedures and named-let loops, the compiler
analyzes which parameters are used exclusively in numeric contexts
(the {\tt analyze-numeric-params} pass).  Parameters meeting three
conditions---every use is numeric-compatible, at least one use
``forces'' a numeric type, and all self-tail-call arguments are
numeric---are unboxed to local {\tt LONGREAL} variables:

\begin{lstlisting}[language=Modula3]
VAR nr_0, nr_1 : LONGREAL;
BEGIN
  nr_0 := NARROW(a_0, SchemeLongReal.T)^;
  nr_1 := NARROW(a_1, SchemeLongReal.T)^;
  LOOP
    (* inner loop uses nr_0, nr_1 directly *)
    result := SchemeLongReal.FromLR(nr_0 + nr_1);
    (* tail-call reassignment: *)
    nr_0 := nr_0 + 1.0d0;
    ...
  END
END
\end{lstlisting}

\noindent This eliminates per-iteration heap allocation and {\tt NARROW}
overhead, yielding 4--5$\times$ speedups on recursive numeric
benchmarks.

%----------------------------------------------------------------------
\section{Impact of a Numeric Tower on Compiled Code}
\label{sec:tower-impact}
%----------------------------------------------------------------------

Introducing exact integers breaks the assumption that all numbers are
{\tt SchemeLongReal.T}.  The compiler must be modified in several
ways.

\subsection{Inline Arithmetic Dispatch}

Currently, inline {\tt +} emits a single {\tt NARROW} to
{\tt SchemeLongReal.T}.  With a numeric tower, the generated code
must dispatch on the runtime types of both operands.  The fastest
approach is a two-level {\tt ISTYPE} cascade:

\begin{lstlisting}[language=Modula3]
(* compiled from (+ a b) with numeric tower *)
IF ISTYPE(a_0, REF INTEGER)
   AND ISTYPE(a_1, REF INTEGER) THEN
  (* fast path: fixnum + fixnum *)
  SchemeInt.Add(a_0, a_1)
ELSE
  (* slow path: coerce to LONGREAL *)
  SchemeLongReal.FromLR(
    SchemeLongReal.FromO(a_0) +
    SchemeLongReal.FromO(a_1))
END
\end{lstlisting}

\noindent The fixnum fast path avoids heap allocation when the result
fits in {\tt INTEGER}; the slow path falls back to the existing
{\tt LONGREAL} code with a generalized {\tt FromO} that handles all
numeric types.

\subsection{Unboxing Under a Tower}

The current unboxing pass assumes all numeric values are {\tt
LONGREAL}.  With a tower, two unboxing strategies are possible:

\begin{enumerate}
\item \textbf{Unbox to {\tt LONGREAL} (conservative).}
  Keep the current scheme.  {\tt NARROW} to {\tt SchemeLongReal.T}
  at loop entry; if the value is actually a fixnum, the {\tt NARROW}
  fails.  The compiler could emit a guard: if all initial values are
  {\tt SchemeLongReal.T}, enter the unboxed loop; otherwise fall back
  to a generic loop.

\item \textbf{Unbox to {\tt INTEGER} (aggressive).}
  When the compiler can prove that loop variables remain in the
  fixnum range (e.g., a counting loop from 0 to $n$), unbox to
  native {\tt INTEGER}.  This is faster than {\tt LONGREAL} for
  pure integer work but requires overflow checking.
\end{enumerate}

The conservative approach preserves existing behavior.  The aggressive
approach is a significant optimization but requires the overflow
detection machinery described in Section~\ref{sec:overflow}.

\subsection{Type Inference Extensions}

The current type inference ({\tt analyze-numeric-params}) classifies
parameters as ``numeric'' or ``generic.''  With a tower, a finer
classification is desirable:

\begin{itemize}
\item \textbf{Integer-only}: parameter used only with integer
  operations ({\tt quotient}, {\tt remainder}, bitwise ops).
\item \textbf{Numeric}: parameter used with general arithmetic
  ({\tt +}, {\tt -}, {\tt *}).
\item \textbf{Generic}: parameter used in non-numeric contexts.
\end{itemize}

Integer-only parameters could be unboxed to {\tt INTEGER} (with
overflow promotion); numeric parameters to {\tt LONGREAL}; generic
parameters remain boxed.

%----------------------------------------------------------------------
\section{Integer Overflow in Modula-3}
\label{sec:overflow}
%----------------------------------------------------------------------

\subsection{What the Language Specification Says}

The Modula-3 language definition~\cite{spwm3,m3defn} is notably
silent on integer overflow semantics.  Integer arithmetic overflow
is \emph{not} explicitly classified as a checked runtime error, an
unchecked runtime error, or undefined behavior.  The specification
states only that programs should use the {\tt Word} interface for
modular arithmetic, implying that ordinary {\tt INTEGER} operations
do \emph{not} wrap---but without specifying what they do instead.

The closest guidance appears in the {\tt FOR} loop specification,
which warns:

\begin{quote}
  If the upper bound of the loop is {\tt LAST(INTEGER)} or
  {\tt LAST(LONGINT)}, it should be rewritten as a {\tt WHILE}
  loop to avoid overflow.
\end{quote}

\noindent This warning to programmers suggests overflow is a hazard
to be avoided, but does not define the consequences.

The runtime error enumeration does include {\tt IntegerOverflow}
(in {\tt RuntimeError.T}), with the message ``Integer result too
large to represent,'' but it is a dead letter: no code path in the
current CM3 runtime raises it.

\subsection{What CM3 Actually Does}

The CM3 C backend ({\tt m3back/src/M3C.m3}) translates {\tt INTEGER}
arithmetic directly to C operators.  The code generator's {\tt add},
{\tt subtract}, and {\tt multiply} procedures all delegate to a
generic {\tt op2} function that emits:

\begin{lstlisting}[language=C]
(left_operand) + (right_operand)
(left_operand) - (right_operand)
(left_operand) * (right_operand)
\end{lstlisting}

\noindent There is no overflow check, no use of
{\tt \_\_builtin\_add\_overflow}, and no post-operation validation.
Signed integer overflow in C is \emph{undefined behavior} per the
C standard (C11~\S6.5/5).

The compiler's constant-folding pass ({\tt m3front/src/exprs/\allowbreak
IntegerExpr.m3}) \emph{does} detect overflow at compile time, using
the arbitrary-precision {\tt TInt} module ({\tt m3middle/src/TInt.m3})
which performs full overflow checking.  But this only affects
compile-time constants; runtime arithmetic is unprotected.

On POSIX platforms, the CM3 runtime ({\tt m3core/src/runtime/\allowbreak
POSIX/\allowbreak RTSignalC.c}) installs signal handlers for {\tt
SIGSEGV}, {\tt SIGBUS}, {\tt SIGPIPE}, etc., but \textbf{not} for
{\tt SIGFPE}.  Integer overflow at runtime is therefore completely
undetected.  On Windows, {\tt EXCEPTION\_INT\_OVERFLOW} is mapped to
{\tt RuntimeError.IntegerOverflow}, but this mapping is moot since
the C backend does not generate instructions that would trigger
the exception.

In practice, on all current targets (AMD64, ARM64, I386), signed
integer overflow wraps in two's complement because the underlying
hardware does so and modern compilers on these platforms generate
wrapping code for simple signed arithmetic---but this is an
accident of implementation, not a guarantee.  An aggressive C
optimizer could in principle exploit the undefined behavior.

\subsection{Implications for a Numeric Tower}

This situation has three consequences for fixnum arithmetic:

\begin{enumerate}
\item \textbf{We cannot rely on overflow trapping.}
  There is no signal, exception, or runtime error when {\tt INTEGER}
  arithmetic overflows.  The program continues with a silently wrong
  value.

\item \textbf{We cannot rely on wrapping.}
  Although wrapping occurs in practice, it is C undefined behavior
  and could change with compiler optimizations.

\item \textbf{Overflow must be detected explicitly before
  performing the operation.}
  This is the only portable approach.
\end{enumerate}

\subsection{Detection Strategies}

\paragraph{Pre-check.}
Before {\tt a + b}, check whether the result would overflow:
\begin{lstlisting}[language=Modula3]
IF (b > 0 AND a > LAST(INTEGER) - b) OR
   (b < 0 AND a < FIRST(INTEGER) - b) THEN
  (* promote to bignum *)
ELSE
  RETURN a + b
END
\end{lstlisting}
This adds two comparisons and a subtraction per operation, roughly
tripling the cost of a fixnum addition.

\paragraph{Post-check via {\tt Word}.}
Perform the addition using {\tt Word.Plus} (unsigned, wrapping---this
\emph{is} well-defined) and check whether the signs indicate overflow:
\begin{lstlisting}[language=Modula3]
WITH sum = Word.Plus(a, b) DO
  IF Word.And(
       Word.Xor(sum, a),
       Word.Xor(sum, b))
     >= Word.Shift(1, BITSIZE(INTEGER)-1) THEN
    (* overflow: promote to bignum *)
  ELSE
    RETURN VAL(sum, INTEGER)
  END
END
\end{lstlisting}
This uses defined-behavior unsigned arithmetic throughout.  The
sign-bit check detects overflow in both directions.  Still adds
overhead, but avoids the branch-heavy pre-check.

\paragraph{Compiler intrinsic via emitted C.}
Since the MScheme compiler emits Modula-3 source (not C), it cannot
directly use GCC/Clang {\tt \_\_builtin\_add\_overflow} intrinsics.
However, a thin Modula-3 wrapper with {\tt <*EXTERNAL*>} pragmas
calling a C helper function is possible:

\begin{lstlisting}[language=C]
int SchemeFixnum__AddCheck(long a, long b, long *result) {
    return __builtin_add_overflow(a, b, result);
}
\end{lstlisting}

\noindent This is the fastest option (single instruction on x86:
{\tt jo} after {\tt add}) but couples the implementation to the
C backend and GCC/Clang.

\paragraph{Practical recommendation.}
Use the {\tt Word}-based post-check for the initial implementation.
It is portable, uses only defined behavior, and has acceptable
overhead.  If profiling shows fixnum arithmetic to be a bottleneck,
add a C helper using {\tt \_\_builtin\_*\_overflow} as an alternative
fast path.

\paragraph{Multiplication.}
Overflow detection for multiplication is harder than for addition.
The pre-check requires a division (expensive) and has edge cases
around {\tt FIRST(INTEGER)}.  The {\tt Word}-based approach requires
a double-width multiply or a comparison against the maximum quotient.
In practice, the simplest correct approach for multiplication is:
check whether both operands fit in 32~bits (a single range check);
if so, the 64-bit product cannot overflow; otherwise, promote both
operands to {\tt Mpz.T} and multiply there.

%----------------------------------------------------------------------
\section{Bootstrap Considerations}
%----------------------------------------------------------------------

The MScheme compiler is self-hosting: it is written in Scheme and
compiled by itself.  Introducing a numeric tower requires a two-phase
bootstrap:

\begin{enumerate}
\item \textbf{Phase~1.}  Compile the new compiler source under the
  old compiler (which uses {\tt LONGREAL}-only arithmetic).  The new
  compiler source must be valid under both the old and new numeric
  semantics.  Since integer arithmetic on doubles is exact for values
  up to $2^{53}$ and the compiler processes only small integers
  (line numbers, indices, string lengths), this is straightforward.

\item \textbf{Phase~2.}  Use the Phase~1 compiler to recompile
  itself.  This produces a compiler that emits numeric-tower-aware
  Modula-3 code.
\end{enumerate}

The bootstrap is eased by the fact that the compiler itself does not
exercise the numeric tower strenuously---it manipulates symbols,
strings, and small integers, all well within {\tt LONGREAL}'s exact
range.

%======================================================================
\chapter{Integer Representation Design}
\label{ch:integer-design}
%======================================================================

This chapter evaluates three concrete designs for exact integer
representation and recommends one.

%----------------------------------------------------------------------
\section{Design A: Fixnum Only ({\tt REF INTEGER})}
\label{sec:fixnum-only}
%----------------------------------------------------------------------

\paragraph{Representation.}
All exact integers are {\tt REF INTEGER}.  Overflow on arithmetic
signals an error (or silently promotes to {\tt LONGREAL} inexact).

\paragraph{Advantages.}
\begin{itemize}
\item Minimal implementation complexity.
\item Fast arithmetic (native {\tt INTEGER} operations).
\item The exact/inexact distinction is established: {\tt REF INTEGER}
  is exact, {\tt REF LONGREAL} is inexact.
\end{itemize}

\paragraph{Disadvantages.}
\begin{itemize}
\item Violates R5RS: overflow on {\tt (+\ MOST-POSITIVE-FIXNUM\ 1)}
  should produce a correct (bignum) result, not an error.
\item 64-bit range is large but not unlimited; edge cases in
  cryptographic or number-theoretic code will fail.
\item {\tt exact->inexact} on large fixnums loses precision
  (inherent in IEEE~754, not a bug).
\end{itemize}

\paragraph{Verdict.}
Acceptable as a transitional step but not R5RS-conformant.

%----------------------------------------------------------------------
\section{Design B: Mpz Only}
\label{sec:mpz-only}
%----------------------------------------------------------------------

\paragraph{Representation.}
All exact integers are {\tt Mpz.T}.  No fixnum representation;
even small integers like $0$ and $1$ are GMP objects.

\paragraph{Advantages.}
\begin{itemize}
\item Correct by construction: no overflow, arbitrary precision.
\item Single code path for all exact integers (no dispatch on
  fixnum vs.\ bignum).
\item GMP internally uses machine-word arithmetic for small values,
  so the constant factor is tolerable.
\end{itemize}

\paragraph{Disadvantages.}
\begin{itemize}
\item Every exact integer requires a GMP allocation (at least 3
  machine words for the {\tt mpz\_t} structure plus the Modula-3
  object header and weak reference for deallocation).
\item Arithmetic on small integers is slower than native {\tt INTEGER}:
  function call overhead, GMP's internal dispatch, memory indirection.
\item External dependency on GMP.
\end{itemize}

\paragraph{Verdict.}
Correct and simple, but unnecessarily slow for the common case
of small integers.

%----------------------------------------------------------------------
\section{Design C: Three-Level (Recommended)}
\label{sec:three-level}
%----------------------------------------------------------------------

\paragraph{Representation.}
\begin{enumerate}
\item \textbf{Fixnum}: {\tt REF INTEGER} for values in
  $[-2^{63}, 2^{63}-1]$.
\item \textbf{Bignum}: {\tt Mpz.T} for values outside fixnum range.
\item \textbf{Inexact}: {\tt REF LONGREAL} (unchanged).
\end{enumerate}

\paragraph{Type dispatch.}
Every numeric operation begins with:
\begin{lstlisting}[language=Modula3]
IF ISTYPE(x, REF INTEGER) THEN
  (* fixnum path *)
ELSIF ISTYPE(x, Mpz.T) THEN
  (* bignum path *)
ELSIF ISTYPE(x, SchemeLongReal.T) THEN
  (* inexact path *)
ELSE
  RAISE E("not a number: " & ...)
END
\end{lstlisting}

\paragraph{Fixnum caching.}
Small fixnums (e.g., $[-256, 256]$) should be pre-allocated in a
global array, analogous to the current {\tt SchemeLongReal.Zero/One}
caching but covering a wider range.  This eliminates allocation for
the most common integer values.

\paragraph{Promotion and demotion.}
\begin{itemize}
\item Fixnum arithmetic that overflows promotes to {\tt Mpz.T}.
\item Mpz results that fit in fixnum range demote to {\tt REF INTEGER}.
  (The {\tt Mpz.ToInteger} function provides the range check.)
\item Mixing exact and inexact (e.g., fixnum + float) produces
  inexact, following R5RS contagion rules.
\end{itemize}

\paragraph{Advantages.}
\begin{itemize}
\item R5RS-conformant: no overflow, no precision loss on exact
  arithmetic.
\item Fast for common cases: small integer arithmetic uses native
  {\tt INTEGER}.
\item GMP available for arbitrary precision when needed.
\item Clear exact/inexact semantics.
\end{itemize}

\paragraph{Disadvantages.}
\begin{itemize}
\item Two-way dispatch on the integer path (fixnum vs.\ bignum).
\item Overflow detection on the fixnum path (Section~\ref{sec:overflow}).
\item More complex code than either single-representation design.
\end{itemize}

\paragraph{Verdict.}
The right trade-off for a production Scheme implementation.  The
three-level design is standard practice in major Scheme and Lisp
implementations (Chez Scheme, Guile, SBCL, Racket) for good reason.

%----------------------------------------------------------------------
\section{Fixnum Caching Detail}
%----------------------------------------------------------------------

A pre-allocated array of {\tt REF INTEGER} values for the range
$[-256, 256]$ (513 entries) eliminates allocation for the vast
majority of integer operations in practice:

\begin{lstlisting}[language=Modula3]
VAR cache : ARRAY [-256..256] OF REF INTEGER;

PROCEDURE FromI(x : INTEGER) : REF INTEGER =
  BEGIN
    IF x >= -256 AND x <= 256 THEN
      RETURN cache[x]
    ELSE
      WITH r = NEW(REF INTEGER) DO r^ := x; RETURN r END
    END
  END FromI;
\end{lstlisting}

\noindent The cache size is a tunable parameter.  Larger caches
reduce allocation at the cost of startup memory; 513 entries
($\approx 4$\,KB on 64-bit) is negligible.

%======================================================================
\chapter{Case Study: Integer Types in the CSP Compiler}
\label{ch:csp}
%======================================================================

The CSP compiler ({\tt cspc}, located at {\tt intel-async/\allowbreak
async-toolkit/\allowbreak m3utils/\allowbreak m3utils/\allowbreak
csp/})~\cite{cspmanual} is an MScheme application that compiles a
hardware description language for asynchronous circuits.  Its
treatment of integers illustrates a more sophisticated approach
to integer types that is instructive for the numeric tower design.

%----------------------------------------------------------------------
\section{Three Kinds of Integers}
%----------------------------------------------------------------------

CSP defines three integer type categories, reflecting the needs
of hardware description:

\begin{description}[style=nextline]
\item[{\tt int} (arbitrary precision).]
  The default integer type.  Expressions are evaluated in arbitrary
  precision---there is no overflow, no wrapping, and no width limit.
  This is the type of intermediate computations.

\item[{\tt int(W)} (unsigned, $W$~bits).]
  Range: $[0, 2^W - 1]$.  Represents a hardware register of width $W$.
  Assignment to an {\tt int(W)} variable implicitly truncates
  (two's complement wrapping).

\item[{\tt sint(W)} (signed, $W$~bits).]
  Range: $[-2^{W-1}, 2^{W-1}-1]$.  Also represents a hardware register,
  but with two's complement signed interpretation.
\end{description}

\noindent The CSP manual states the design principle explicitly:

\begin{quote}
  Expressions use arbitrary precision arithmetic, but variables can
  either have arbitrary precision OR specified precision, in which
  case the programmer-designer specifies the precise bit width of
  the variable (and whether it is signed or not).
\end{quote}

This separation---\emph{expressions} are arbitrary-precision,
\emph{storage} may truncate---is directly analogous to the Scheme
numeric tower's treatment of exact integers (arbitrary precision
in computation, with {\tt exact->inexact} as the explicit lossy
conversion).

%----------------------------------------------------------------------
\section{Interval Arithmetic for Range Tracking}
%----------------------------------------------------------------------

The CSP compiler performs static range analysis using interval
arithmetic over an extended number system ({\tt xnum}).  An
{\tt xnum} is either a {\tt BigInt.T} (arbitrary precision),
{\tt +inf}, {\tt -inf}, or {\tt nan}.  A \emph{range} is a pair
$[\mathit{min}, \mathit{max}]$ of xnums.

Ranges propagate through all arithmetic operations via the
\emph{cartesian product principle}: for a binary operation
$\mathop{op}$ on ranges $[a_{\min}, a_{\max}]$ and
$[b_{\min}, b_{\max}]$, compute all four corner combinations
$\mathop{op}(a_i, b_j)$ and return the range
$[\min(\text{results}), \max(\text{results})]$.

From {\tt csp/src/bigint.scm}:

\begin{lstlisting}[language=Scheme]
(define (make-simple-range-binop op)
  (lambda (a b)
    (let* ((all-pairs   (cartesian-product a b))
           (all-results (map eval
             (map (lambda (x) (cons op x))
                  all-pairs)))
           (min-res (apply xnum-min all-results))
           (max-res (apply xnum-max all-results)))
      (list min-res max-res))))

(define range-+ (make-simple-range-binop xnum-+))
(define range-* (make-simple-range-binop xnum-*))
\end{lstlisting}

\noindent Division requires special handling: the range is split at
zero to handle sign changes, and division by zero yields
$+\infty$ or {\tt nan}.

%----------------------------------------------------------------------
\section{Bitwise Operations on Ranges}
%----------------------------------------------------------------------

The CSP compiler also performs range analysis for bitwise
operations---a problem significantly harder than arithmetic ranges,
since bitwise operations do not preserve the interval structure in
the same way.  The approach uses two helper functions:

\begin{itemize}
\item {\tt xnum-setallbits}: given $x$, returns the value with all
  bits set from bit~0 to the MSB of $x$ (i.e., $2^{\lceil\log_2(x+1)\rceil}-1$).
  This is the \emph{maximum} value achievable by OR-ing arbitrary
  bit patterns up to the width of $x$.

\item {\tt xnum-clearallbits}: the complement operation, returning
  the \emph{minimum} value achievable by AND-ing.
\end{itemize}

For example, the OR range of two non-negative ranges is:

\begin{lstlisting}[language=Scheme]
(define (range-| a b)
  (let* ((rmin (xnum-min (range-min a) (range-min b)))
         (rmax (xnum-| (xnum-setallbits (range-max a))
                       (xnum-setallbits (range-max b)))))
    (make-range rmin rmax)))
\end{lstlisting}

\noindent XOR ranges require considering sign interactions and are
the most complex, handling positive and negative sub-ranges
separately.

The implementation comment notes:

\begin{quote}
  Interval arithmetic on bitwise operators is not entirely
  straightforward, and some of the intervals used here may not be
  tight.  It is however believed they are all sound as coded.
\end{quote}

\noindent ``Sound but not tight'' is the practical standard: ranges
may be conservatively wider than necessary, but never exclude valid
values.

%----------------------------------------------------------------------
\section{Relevance to MScheme}
%----------------------------------------------------------------------

The CSP compiler's approach suggests two lessons for MScheme's
numeric tower:

\begin{enumerate}
\item \textbf{Arbitrary-precision expressions with bounded storage}
  is a proven pattern.  The Scheme equivalent is: exact integer
  arithmetic is arbitrary-precision; conversion to fixnum-width
  storage (e.g., for Modula-3 interop via {\tt SchemeLongReal.Int})
  is an explicit, checked operation.

\item \textbf{Bitwise operations on bignums are well-defined} in the
  two's complement model.  The CSP compiler's range analysis for
  bitwise operations, while approximate, demonstrates that the
  operations themselves are straightforward---the subtlety is only in
  static analysis, which MScheme's dynamic types sidestep entirely.
  GMP's bitwise operations ({\tt mpz\_and}, {\tt mpz\_ior}, etc.)
  already implement two's complement semantics for negative operands.
\end{enumerate}

%----------------------------------------------------------------------
\section{What a Numeric Tower Would Simplify in {\tt cspc}}
\label{sec:cspc-simplification}
%----------------------------------------------------------------------

The CSP compiler provides a concrete case study of the cost imposed
by MScheme's lack of exact integers.  Because all MScheme numbers
are {\tt LONGREAL}, {\tt cspc} maintains its own arbitrary-precision
integer layer in {\tt bigint.scm} (1\,009~lines), built on top of
Modula-3's {\tt BigInt.T} via the Scheme/Modula-3 stub interface.

This file has four distinct layers:

\begin{description}[style=nextline]
\item[Layer~1: {\tt BigInt.T} wrappers ($\sim\!85$~lines).]
  Functions like {\tt force-bigint} (convert a Scheme number to
  {\tt BigInt.T}), {\tt bigint?} (type test), and {\tt make-big-op}
  (template for wrapping {\tt BigInt} operations as Scheme
  procedures).  These exist \emph{solely} because MScheme's native
  numbers cannot represent exact integers.
  \textbf{With native bignums: delete entirely.}

\item[Layer~2: Extended numbers / {\tt xnum} ($\sim\!280$~lines).]
  An ``extended number'' is either a {\tt BigInt.T}, {\tt +inf},
  {\tt -inf}, or {\tt nan}.  Each {\tt xnum} arithmetic function
  ({\tt xnum-+}, {\tt xnum-*}, etc.) dispatches on these four
  cases: infinity/NaN handling is algorithmic necessity, but the
  {\tt BigInt.T} branch is pure wrapper code
  (e.g., {\tt (big+ a b)} instead of {\tt (+ a b)}).
  \textbf{With native bignums: $\sim\!120$~lines of wrapper calls
  become direct Scheme arithmetic; $\sim\!130$~lines of
  $\pm\infty$/NaN handling remain unchanged.}

\item[Layer~3: Range/interval system ($\sim\!380$~lines).]
  The cartesian-product range operations ({\tt range-+},
  {\tt range-*}, {\tt range-|}, etc.) are already abstracted
  through the {\tt xnum} layer.
  \textbf{With native bignums: no changes.}

\item[Layer~4: Utilities and tests ($\sim\!160$~lines).]
  Evaluation helpers ({\tt xnum-eval}) and test code.
  \textbf{With native bignums: can be simplified or deleted.}
\end{description}

\noindent The quantitative impact is summarized in
Table~\ref{tab:cspc-impact}.

\begin{table}[htbp]
\centering
\begin{tabular}{lrrr}
\toprule
\textbf{Category} & \textbf{Current} & \textbf{With bignums} & \textbf{Change} \\
\midrule
{\tt BigInt.T} wrappers         &  85 &   0 & $-100\%$ \\
{\tt bigint?} type dispatch     &  30 &   0 & $-100\%$ \\
{\tt xnum-*} wrapper calls      & 120 &  $\sim\!0$ & $-100\%$ \\
$\pm\infty$/NaN handling        & 130 & 130 & $0\%$ \\
Range/interval logic            & 380 & 380 & $0\%$ \\
Utilities/tests                 & 160 & 160 & $0\%$ \\
\midrule
\textbf{Total}                  & \textbf{1\,009} & \textbf{$\sim\!770$} & $\mathbf{-24\%}$ \\
\bottomrule
\end{tabular}
\caption{Impact of native exact integers on {\tt cspc}'s
  {\tt bigint.scm}.  The 380~lines of range logic are untouched
  because they are already abstracted through the {\tt xnum} layer.}
\label{tab:cspc-impact}
\end{table}

Beyond line counts, the deeper benefit is the elimination of an
entire abstraction layer.  Currently, {\tt cspc} must maintain a
parallel number system: every arithmetic expression goes through
{\tt force-bigint} to convert from Scheme's {\tt LONGREAL} to
{\tt BigInt.T}, performs the operation via Modula-3 stubs, and
converts back.  There are 21~calls to {\tt BigInt.New} ({\tt
LONGREAL} $\to$ {\tt BigInt.T}) and 41~{\tt bigint?} type checks
scattered through the code.  With native bignums, ordinary Scheme
arithmetic ({\tt +}, {\tt *}, {\tt expt}, etc.) would be exact by
default, and the entire conversion apparatus vanishes.

The {\tt xnum} layer would still be needed for $\pm\infty$ and
NaN semantics, but its finite-integer branches would reduce from
multi-line wrapper calls to direct Scheme operations---a one-line
change per function.

Two categories of {\tt BigInt} calls would \emph{remain} even with
native bignums: {\tt BigInt.ToInteger} (13~calls, extracting machine
integers for Modula-3 code generation) and {\tt BigInt.Format}
(18~calls, converting to string for code output).  These represent
the genuine Modula-3 interface boundary, not a workaround for
missing Scheme functionality.  With a numeric tower, these would
become calls to {\tt number->string} or the {\tt SchemeLongReal.Int}
bridge.

%======================================================================
\chapter{Exotic Representations}
\label{ch:exotic}
%======================================================================

This chapter briefly surveys alternative numeric representations
that have been proposed in the literature but are not recommended
for MScheme.

%----------------------------------------------------------------------
\section{Continued Fractions}
%----------------------------------------------------------------------

Continued fractions represent real numbers as:
\[
  a_0 + \cfrac{1}{a_1 + \cfrac{1}{a_2 + \cfrac{1}{\ddots}}}
\]
Rational numbers have finite continued-fraction representations;
irrationals have infinite ones.  Arithmetic on continued fractions
(Gosper's algorithm~\cite{gosper}) processes terms lazily, enabling
exact real arithmetic for algebraic expressions.

\paragraph{Why not for MScheme.}
Continued-fraction arithmetic is complex to implement and slow
in practice.  Basic operations like addition require maintaining
a matrix of four bignums per partial quotient; comparison can require
unbounded lookahead.  The result quality (exact for rationals, lazy
for irrationals) does not justify the implementation complexity for
a Scheme primarily used as a scripting language.  Standard
arbitrary-precision rationals (Section~\ref{sec:exact-integers})
provide exact arithmetic for all rational values with much simpler
code.

%----------------------------------------------------------------------
\section{Interval Arithmetic}
%----------------------------------------------------------------------

Interval arithmetic represents a real number as a pair $[a, b]$
where the true value lies in the interval.  Operations propagate
intervals, maintaining rigorous error bounds.

\paragraph{Why not for MScheme.}
Interval arithmetic addresses a different problem (numerical error
tracking) than the Scheme numeric tower (exact/inexact distinction).
It could be implemented as a \emph{library} type rather than a core
numeric type.  The overhead of maintaining two bounds per number is
significant, and standard Scheme code does not expect interval
semantics.

%----------------------------------------------------------------------
\section{$p$-adic Numbers}
%----------------------------------------------------------------------

The $p$-adic numbers $\mathbb{Q}_p$ provide an alternative
completion of the rationals, where closeness is measured by
divisibility by powers of a prime $p$.  They appear in number
theory and algebraic geometry.

\paragraph{Why not for MScheme.}
$p$-adic numbers are a specialized mathematical tool with essentially
no presence in the Scheme standard or common Scheme usage patterns.
If needed, they could be implemented as a library type.

%======================================================================
\chapter{Verification and Testing}
\label{ch:testing}
%======================================================================

%----------------------------------------------------------------------
\section{The Impracticality of Full Formal Verification}
%----------------------------------------------------------------------

Formally verifying a numeric tower implementation---proving that
every operation produces correct results for all inputs---is
attractive in principle but impractical here.

\begin{itemize}
\item The implementation spans two languages (Scheme compiler output
  in Modula-3, compiled by CM3 to C, compiled by GCC/Clang to native
  code) plus an external library (GMP).  Verification would need to
  cross all these boundaries.

\item GMP itself is not formally verified.  Its correctness rests
  on decades of testing and widespread use, not proof.  A formal
  verification of MScheme's tower would need to treat GMP as an
  axiomatically correct oracle.

\item The Modula-3 compiler and runtime are not formally verified.
  {\tt NARROW}, garbage collection, and weak references are all
  trusted.

\item The effort required for formal verification (in, e.g.,
  Coq or Isabelle) would dwarf the implementation effort by orders
  of magnitude, and the resulting proofs would be fragile against
  implementation changes.
\end{itemize}

%----------------------------------------------------------------------
\section{Property-Based Testing}
%----------------------------------------------------------------------

The practical alternative is thorough testing, combining traditional
unit tests with property-based (randomized) testing.

\paragraph{Algebraic properties.}
Test that arithmetic satisfies algebraic laws across all
representation boundaries:

\begin{itemize}
\item Commutativity: $a + b = b + a$, $a \times b = b \times a$.
\item Associativity: $(a + b) + c = a + (b + c)$.
\item Identity: $a + 0 = a$, $a \times 1 = a$.
\item Inverse: $a + (-a) = 0$, $a \times (1/a) = 1$ (for $a \neq 0$).
\item Distributivity: $a \times (b + c) = a \times b + a \times c$.
\end{itemize}

\noindent These must hold for all combinations of fixnum, bignum,
rational, and inexact operands (with appropriate tolerance for
inexact results).

\paragraph{Boundary testing.}
Focus on representation boundaries:

\begin{itemize}
\item Fixnum limits: {\tt FIRST(INTEGER)}, {\tt LAST(INTEGER)},
  and values $\pm 1$ from these.
\item The fixnum/bignum promotion boundary: operations whose inputs
  are fixnums but whose results require bignums.
\item The bignum/fixnum demotion boundary: bignum operations that
  produce fixnum-range results.
\item The exact/inexact boundary: mixing exact and inexact operands.
\item The IEEE~754 exact-integer limit ($2^{53}$): values near
  this boundary that would lose precision when converted to {\tt
  LONGREAL}.
\end{itemize}

\paragraph{Conversion round-trips.}
\begin{itemize}
\item {\tt (string->number (number->string $x$)) = $x$} for exact $x$.
\item {\tt (inexact->exact (exact->inexact $x$)) = $x$} when $x$ is
  a small exact integer.
\item {\tt (exact->inexact $x$)} preserves value when $|x| \leq 2^{53}$.
\end{itemize}

\paragraph{Cross-validation.}
Run R5RS test suites (e.g., the R7RS test suite from Chibi-Scheme,
which includes extensive numeric tests) and compare results against
a reference implementation like Chez Scheme or Guile.

%----------------------------------------------------------------------
\section{Test Infrastructure}
%----------------------------------------------------------------------

MScheme already has a test framework ({\tt test-r4rs-gaps.scm}) that
can be extended.  Suggested additions:

\begin{enumerate}
\item An ``exact integer'' test suite exercising fixnum/bignum
  boundaries.
\item A ``rational arithmetic'' test suite with known-answer tests
  (e.g., $1/3 + 1/6 = 1/2$).
\item A ``mixed exactness'' test suite verifying contagion rules.
\item A property-based test harness that generates random operands
  across representation boundaries.
\end{enumerate}

%======================================================================
\chapter{Conclusions and Recommendations}
\label{ch:conclusions}
%======================================================================

%----------------------------------------------------------------------
\section{Summary of Findings}
%----------------------------------------------------------------------

\begin{enumerate}
\item MScheme's single-{\tt LONGREAL} representation is adequate for
  R4RS but insufficient for R5RS conformance.  The most critical gap
  is the collapsed exact/inexact distinction.

\item The three-level integer representation (Section~\ref{sec:three-level})---fixnum
  ({\tt REF INTEGER}), bignum ({\tt Mpz.T}), inexact ({\tt REF
  LONGREAL})---is the standard design for a reason: it is correct,
  fast for common cases, and well-understood.

\item The {\tt Mpz} GMP bindings already exist with over 200
  operations, are typed as {\tt REFANY}, use auto-generated
  bindings (easily updated), and include WeakRef-based memory
  management.  GMP's {\tt mpq} rational functions are present in
  the generator database but not yet exposed; creating {\tt Mpq.T}
  bindings is straightforward.  MPFR bindings also exist for
  arbitrary-precision floats.

\item The {\tt SchemeLongReal.FromO} bridge function is the
  critical backward-compatibility mechanism.  Extending it to
  handle exact integer types preserves all existing Modula-3
  interoperability.

\item \textbf{Integer overflow in CM3 is undefined behavior.}
  The Modula-3 specification does not classify integer overflow;
  the CM3 C backend emits plain C signed arithmetic with no overflow
  checks; the runtime does not catch {\tt SIGFPE}.  Fixnum arithmetic
  \emph{must} use explicit pre-checks or {\tt Word}-based post-checks.
  The overflow detection cost is the primary performance concern for
  the fixnum path.

\item The compiler's unboxing optimization can initially be left
  unchanged (unboxing to {\tt LONGREAL}).  A second phase could
  add {\tt INTEGER} unboxing for counting loops with overflow
  promotion.

\item {\tt IntForRat} has a documented silent-overflow bug and
  must not be used for exact rationals.  Use {\tt BigRational}
  (with {\tt BigInt}) or create {\tt MpzRational}.

\item The CSP compiler demonstrates a proven pattern of
  arbitrary-precision expressions with bounded-width storage, and
  shows that bitwise operations on intervals are tractable.
  GMP natively supports two's complement bitwise operations on
  negative bignums, directly enabling SRFI-151.
\end{enumerate}

%----------------------------------------------------------------------
\section{Recommended Implementation Order}
%----------------------------------------------------------------------

\begin{enumerate}
\item \textbf{Phase~1: Exact integers (fixnum + bignum).}
  Introduce {\tt REF INTEGER} and {\tt Mpz.T} as exact integer
  types.  Modify {\tt IsExact}, {\tt FromO}, {\tt Int}, and the
  primitive dispatch in {\tt SchemePrimitive.m3}.  Update the
  compiler's inline arithmetic to emit fixnum fast paths.  This
  alone achieves R5RS integer compliance.

\item \textbf{Phase~2: Exact rationals.}
  Add {\tt BigRational.T} (or {\tt MpzRational.T}) as the exact
  rational type.  Implement {\tt numerator}, {\tt denominator},
  {\tt rationalize}, and update {\tt /} to return exact rationals
  when dividing exact integers with a non-zero remainder.

\item \textbf{Phase~3: Complex numbers.}
  Add {\tt SchemeComplex.T} as described in the representations
  chapter.  Implement {\tt make-rectangular}, {\tt make-polar},
  {\tt real-part}, {\tt imag-part}, {\tt magnitude}, {\tt angle}.

\item \textbf{Phase~4: Compiler optimizations.}
  Add {\tt INTEGER}-unboxed loops, fixnum fast-path inlining, and
  SRFI-143/151 bitwise operations.
\end{enumerate}

Each phase is self-contained and produces a working system.  Phase~1
delivers the most value per effort; Phases~2 and~3 can be deferred
indefinitely if the application domain does not require them.

%----------------------------------------------------------------------
\section{Effort Estimate by Phase}
%----------------------------------------------------------------------

\begin{table}[htbp]
\centering
\begin{tabular}{llc}
\toprule
\textbf{Phase} & \textbf{Scope} & \textbf{New/Modified M3 Lines} \\
\midrule
1. Exact integers   & Interpreter + compiler + sstubgen & $\sim$500--800 \\
2. Exact rationals  & Interpreter + reader/printer      & $\sim$300--500 \\
3. Complex numbers  & Interpreter + reader/printer      & $\sim$300--400 \\
4. Compiler opts    & Compiler only                     & $\sim$200--400 \\
\bottomrule
\end{tabular}
\caption{Estimated effort by phase.  Phase~1 is the largest because
  it touches the most subsystems (interpreter dispatch, compiler code
  generation, sstubgen stubs, reader/printer).  Subsequent phases are
  more localized.}
\label{tab:effort}
\end{table}

\begin{thebibliography}{99}

\bibitem{r4rs} William Clinger and Jonathan Rees, eds.
Revised$^4$ Report on the Algorithmic Language Scheme.  1991.

\bibitem{r5rs} Richard Kelsey, William Clinger, and Jonathan Rees, eds.
Revised$^5$ Report on the Algorithmic Language Scheme.  1998.

\bibitem{r7rs} Alex Shinn, John Cowan, and Arthur A.\ Gleckler, eds.
Revised$^7$ Report on the Algorithmic Language Scheme.  2013.

\bibitem{spwm3} Greg~Nelson, ed.  {\it Systems Programming with Modula-3.}
Englewood Cliffs, N.J.:\ Prentice-Hall, 1991.

\bibitem{jscheme} Peter Norvig.  JScheme.
{\tt https://norvig.com/jscheme.html}

\bibitem{r4rsconf} Mika Nystr\"om.  MScheme R4RS Conformance Report.
March 2026.

\bibitem{m3defn} Luca Cardelli, James Donahue, Lucille Glassman,
Mick Jordan, Bill Kalsow, and Greg Nelson.
Modula-3 Language Definition.
In~\cite{spwm3}, Chapter~2.

\bibitem{cspmanual} Rajit Manohar.
The CSP Language Reference Manual.
Intel Labs, 2024.

\bibitem{gosper} Bill Gosper.  Continued Fraction Arithmetic.
MIT AI Lab, 1972.  HAKMEM Item~101B.

\end{thebibliography}

\end{document}
